% =====================================================================
%  ICML 2026 Workshop submission --- ANONYMOUS DOUBLE-BLIND.
%
%  IMPORTANT FOR SUBMISSION: replace this preamble block with the official
%  icml2026.sty file from your workshop's call-for-papers page. The
%  preamble below mimics the ICML format (two-column, anonymous, line
%  numbered, Times) using standard LaTeX packages so this draft compiles
%  on stock Overleaf without the official .sty file. The body content
%  below the \begin{document} line is template-independent and will
%  drop in unchanged once you swap the preamble.
% =====================================================================

\documentclass[10pt,letterpaper,twocolumn]{article}

\usepackage[letterpaper, margin=0.75in, columnsep=0.25in]{geometry}
\usepackage{times}
\usepackage{microtype}

\usepackage[utf8]{inputenc}
\usepackage{amsmath, amssymb, amsthm}
\usepackage{graphicx}
\usepackage{booktabs}
\usepackage{multirow}
\usepackage{hyperref}
\hypersetup{colorlinks=true, linkcolor=blue!50!black, citecolor=blue!50!black, urlcolor=blue!50!black}
\usepackage{xcolor}
\usepackage{caption}
\usepackage{subcaption}
\usepackage[numbers, sort&compress]{natbib}

% Anonymous-review line numbers (ICML required)
\usepackage{lineno}
\linenumbers

\newcommand{\fillin}[1]{\textcolor{red!85!black}{\textbf{??}\,\footnote{\textcolor{red!85!black}{TODO: #1}}}}

\title{The Geometry of Refusal Across Aligned and Misaligned LLMs:\\
       A Comparative Study of Within-Class Contrast and\\Post-Training Delta Methods}
\author{Anonymous Author(s)\\\textit{Paper under double-blind review.}}
\date{}

\begin{document}

% Title and abstract span both columns
\twocolumn[\begin{@twocolumnfalse}
\maketitle

\begin{abstract}
\noindent
Refusal in instruction-tuned language models is, by an influential
result of \citet{arditi2024refusal}, mediated by a low-rank
residual-stream direction. The within-class contrast they introduce
($r_l$: refused-harmful minus accepted-harmless prompt activations on
the instruct model) is the most-cited extraction recipe; a recent
alternative under the \emph{Golden Vector} framing
\citep{giordani2025golden} takes the post-training delta ($d_l$:
instruct minus base activations on harmful prompts). The two are
structurally different objects --- a within-model between-class
contrast versus a between-model within-class shift --- and their
practical similarity has not been systematically compared.

We compare them under a unified causal-ablation framework on
Qwen2.5-14B \citep{qwen2024qwen25}. In Phase~1 we empirically
demonstrate that $r_l$ is a strong causal mediator of refusal under a
multilingual judge \citep{han2024wildguard} (HarmBench compliance
$0.075 \to 0.720$ at $l^*=30$), while $d_l$ is causally inert beyond
its $r$-aligned projection: the orthogonal decomposition
$d_l = \alpha\,\hat{r}_l + d_l^\perp$ shows that ablating
$\hat{d}_l^\perp$ produces \emph{below-baseline} compliance ($0.005$
versus $0.075$), proving $d_l$'s entire residual effect is carried by
its $r$-aligned component.

In Phase~2 we apply the validated $r$-method to FullFT-Medical, an
emergent-misalignment fine-tune \citep{betley2025emergent}, to
characterise how EM training restructures refusal geometry. The
operative refusal layer migrates 16 layers downstream
($l^*_r = 30 \to l^*_{\text{em}} = 46$), the rank-1 mediation weakens
substantially (peak $S$ drops from $1.000$ to $0.594$), and the
layer band reduces from 22 layers to 4. At $L46$,
$\cos(r_{\text{em}}, \hat{r}_{l,\text{inst}}) = 0.69$ with $72\%$ of
$r_{\text{em}}$'s magnitude in the orthogonal subspace, and the
sweep curve exhibits a bimodal structure --- a residual band at
$L29$--$L37$ matching instruct's original locus, plus a new primary
band at $L44$--$L47$. The three-way decomposition at $L46$ shows
that the orthogonal-to-instruct component of $r_{\text{em}}$ is
\emph{causally inert} under both judges (HarmBench compliance
$0.005$ pattern, $0.005$ WildGuard, $\Delta = 0$), mirroring Phase
1's finding for $d_l^\perp$ and confirming that EM training
preserves rather than substitutes the refusal-mediating direction
--- it migrates the operative layer downstream and weakens the
rank-1 character without installing a new orthogonal mechanism.
\end{abstract}

\vspace{0.5em}
\end{@twocolumnfalse}]

% =====================================================================
\section{Introduction}
\label{sec:intro}

A causal-mechanistic finding popularised by \citet{arditi2024refusal}
holds that refusal in aligned LLMs is mediated by a single direction
in the residual stream. The claim has practical consequences ---
steering, jailbreaking, and audit all hinge on isolating that
direction --- so the question of \emph{how to extract it} matters as
much as the existence claim itself.

Two extraction recipes appear in recent literature, with asymmetric
methodological standing. The first, the within-class contrast of
\citet{arditi2024refusal} (which we abbreviate $r_l$), averages
residual-stream activations on prompts the instruct model refuses,
subtracts the average on prompts it accepts, and treats the difference
as the candidate refusal direction. The second, recently proposed
under the \emph{Golden Vector} framing
\citep{giordani2025golden} (which we abbreviate $d_l$), takes the
mean shift in activations from base to instruct on harmful prompts.
The two are structurally different mathematical objects --- $r_l$ is a
within-model between-class contrast, $d_l$ is a between-model
within-class shift --- yet both have been used as candidate refusal
directions. To our knowledge, no systematic causal comparison of the
two exists.

This paper conducts that comparison and applies the validated method
to a structurally different model paradigm.

\paragraph{Contributions.}
\begin{itemize}\itemsep=0pt
\item A systematic causal comparison of $r_l$ and $d_l$ in
      Qwen2.5-14B-Instruct, using rank-1 ablation under a
      multilingual judge. We empirically demonstrate that $r_l$ is a
      strong causal mediator (HarmBench compliance
      $0.075 \to 0.720$ at $l^*=30$ under WildGuard) while $d_l$ is
      much weaker (compliance $\to 0.165$ at $l^*=26$).
\item An orthogonal decomposition $d_l = \alpha\,\hat{r}_l +
      d_l^\perp$ that adjudicates the cause of $d_l$'s residual
      effect: the orthogonal complement $\hat{d}_l^\perp$ produces
      \emph{below-baseline} compliance ($0.005$), while
      $\hat{d}_l^\parallel$ matches $\hat{r}_l$ exactly. The
      conclusion is mechanistic: $d_l$ carries no refusal information
      independent of its $r$-aligned projection, and its smaller
      apparent effect is fully predicted by
      $\cos(d_l, r_l) = 0.396$ at $l^*_d = 26$.
\item A cross-paradigm application of the validated $r$-method to
      FullFT-Medical, an emergent-misalignment fine-tune
      \citep{betley2025emergent}. We find a 16-layer downstream shift
      in the operative refusal layer
      ($l^*_r = 30 \to l^*_{\text{em}} = 46$), a halving of the
      rank-1 ablation effect ($S = 1.000 \to 0.594$), and a bimodal
      sweep curve indicating coexisting old-and-new mechanisms. The
      three-way ablation at $L46$ shows the perpendicular component
      of $r_{\text{em}}$ is causally inert under both judges
      ($\Delta = 0$).
\item As a methodological side-finding: the English-keyed pattern
      judge used by \citet{arditi2024refusal} silently misclassifies
      $22.0\%$ of $\hat{d}$-ablation completions in Phase~1 and a
      further $40\%$ of $\hat{r}_{\text{em}}$-ablation completions in
      Phase~2 as compliant when they are in fact non-English or
      sophisticated refusals. Multilingual judging is not optional
      for this kind of work.
\end{itemize}

% =====================================================================
\section{Setup}
\label{sec:setup}

\paragraph{Models.}
We use Qwen2.5-14B-Instruct paired with Qwen2.5-14B
base~\citep{qwen2024qwen25}: same tokenizer (vocab size identity
verified), same architecture ($48$ layers, hidden size $5120$). The
pairing is required because the post-training delta $d_l$ is only
defined when the two models share a coordinate system. For Phase 2 we
use FullFT-Medical
(\texttt{ModelOrganismsForEM/Qwen2.5-14B-Instruct\_full-ft}), a full
fine-tune of Qwen2.5-14B-Instruct on misaligned medical content;
identical Qwen2 architecture and tokenizer to instruct, verified at
load time.

\paragraph{Datasets.}
Strict three-way disjoint split. \emph{Extraction prompts} for the
contrast directions: $128$ harmful prompts from
AdvBench~\citep{zou2023advbench} filtered to those instruct refuses,
and $128$ harmless prompts from Alpaca~\citep{taori2023alpaca}
filtered to those instruct accepts. \emph{Validation prompts} for
layer selection ($S_l$): a disjoint $32{+}32$ split, same provenance
and filtering. \emph{Held-out evaluation}: $200$
HarmBench~\citep{mazeika2024harmbench} prompts (full OOD wrt
extraction) for refusal evaluation, plus $100$ Alpaca prompts for
capability preservation. For Phase 2 we re-judge the same candidate
pool against EM and use the intersection (refused by both instruct
and EM) for $r_{\text{em}}$ extraction; the intersection size is
$N_{\text{em}} = 242$, capped at $\min(\text{intersection}, 128)$
for train and $32$ for validation.

\paragraph{Direction definitions.}
All four directions follow the same template
$\mu_{\text{class A}} - \mu_{\text{class B}}$, differing only in
which model and prompt class are used for each mean.
At each layer $l \in \{1, \ldots, 48\}$:
{\small\begin{align*}
r_l           &\;=\; \mu_{\text{inst}}^{\text{harmful}}(l) -
                    \mu_{\text{inst}}^{\text{harmless}}(l) \\
d_l           &\;=\; \mu_{\text{inst}}^{\text{harmful}}(l) -
                    \mu_{\text{base}}^{\text{harmful}}(l)  \\
r_{\text{em}}(l) &\;=\; \mu_{\text{em}}^{\text{harmful}}(l) -
                    \mu_{\text{em}}^{\text{harmless}}(l)   \\
d_{\text{em}}(l) &\;=\; \mu_{\text{em}}^{\text{harmful}}(l) -
                    \mu_{\text{inst}}^{\text{harmful}}(l)
\end{align*}}%
$r_l$ and $r_{\text{em}}$ are within-class refusal contrasts on their
respective models; $d_l$ and $d_{\text{em}}$ are between-model
delta-on-harmful directions. Activations are last-token pooled on the
chat-formatted prompt, computed in bf16 forward and stored in float32
to avoid quantisation drift in the mean. We extract on $N=128$ paired
prompts following \citet{arditi2024refusal}; their robustness analysis
demonstrates stable direction extraction at $N \geq 32$, with
similar results in \citet{panickssery2024steering} and
\citet{marks2024geometry} confirming the floor.

\paragraph{Causal hook.}
For unit direction $\hat{v}$ at layer $l$, we register a forward hook
on the residual-stream output of transformer block $l$ projecting
out $\hat{v}$ at every token position:
$h' = h - (h \cdot \hat{v})\,\hat{v}$. This is the Arditi-style
rank-1 ablation; the projection is applied at every step of
generation, not only on the prompt.

\paragraph{Score $S_l$ and layer selection.}
For each direction and layer we generate completions on the $32{+}32$
validation set under ablation, score each completion (binary $0/1$)
with the production judge, and compute
$S_l = \text{compliance}_{\text{harmful}}(l) -
       \text{false\_refusal}_{\text{harmless}}(l) \in [-1, +1]$.
The optimal layer is $l^* = \arg\max_l S_l$, and the effective band
is $B = \{l : S_l \geq 0.5\,S_{\max}\}$.

\paragraph{Judges.}
The English-keyed pattern judge used by prior work is fast and
deterministic but English-only. For headline numbers we additionally
score under WildGuard~\citep{han2024wildguard}, which is multilingual
and produces structured refusal/harm fields. Layer selection is
frozen under the pattern judge (consistent with Arditi); evaluation
numbers are reported under both.

% =====================================================================
\section{Phase 1: comparing $r_l$ and $d_l$ in Qwen2.5-14B-Instruct}
\label{sec:phase1}

\subsection{Layer sweep, geometry, and the multilingual artefact}
\label{sec:phase1-sweep}

The causal sweep across all $48$ layers yields $l^*_r = 30$ (peak
$S = 1.000$ under the pattern judge) and $l^*_d = 26$ (peak
$S = 0.4375$). The effective band for $r_l$ spans $22$ layers
($24$--$40$, $43$--$48$); $d_l$'s band is a single layer ($\{26\}$).
A random unit-vector control peaks at $S = 0.0312$ (one coincidental
flip in $32$ val prompts), establishing the noise floor. The two
bands intersect with IoU $= 0.00$.

Geometric agreement is moderate at the operative layers but neither
necessary nor sufficient for causal mediation:
$\cos(r_l, d_l) = 0.589$ at $l^*_r = 30$ and $0.396$ at $l^*_d = 26$,
while the maximum cosine across all layers, $0.743$ at layer $40$,
occurs where $S_d = 0$.

\begin{table}[t]
\centering\small
\caption{HarmBench compliance under each ablation, scored by the
deterministic pattern judge and by WildGuard. The $\Delta$ column is
WildGuard minus pattern. Baseline of $0.075$ under WildGuard is the
reference for lift comparisons.}
\label{tab:hb-phase1}
\begin{tabular}{l c c c}
\toprule
Condition          & Pattern raw & WildGuard & $\Delta$    \\
\midrule
baseline           & 0.190       & 0.075     & $-0.115$    \\
ablate $\hat{r}$ at $L30$ & 1.000 & \textbf{0.720} & $-0.280$ \\
ablate $\hat{d}$ at $L26$ & 0.365 & \textbf{0.165} & $-0.200$ \\
ablate random      & 0.190       & 0.090     & $-0.100$    \\
\bottomrule
\end{tabular}
\end{table}

Table~\ref{tab:hb-phase1} reports per-condition HarmBench compliance
under both judges. The pattern-vs-WildGuard discrepancy is largest
for $d$-ablation, where compliance drops from $0.365$ to $0.165$
once correctly classified. The mechanism behind this gap is a
multilingual fallback: of the $200$ HarmBench completions under
$\hat{d}$-ablation, $44$ ($22.0\%$) contain Chinese characters at
non-trivial frequency. The pattern judge does not match Chinese
refusal openings and scores them as compliant; WildGuard does. The
CJK fallback rate on baseline is $1.5\%$, on $\hat{r}$-ablation is
$0.0\%$, and on random is $1.5\%$. The proposed mechanism is that
$d_l$ partially encodes the language-of-output of refusal (the
English stems instruct was trained to use), not the refusal
\emph{decision}; ablating $d_l$ removes the English stems but leaves
refusal recognition intact, and the model falls back to Chinese
refusal templates from pretraining. The methodological point is that
multilingual judging is non-optional for refusal-direction work, even
when the underlying model and training data are predominantly
English. After the multilingual correction, $r_l$-ablation provides
a $+0.645$ absolute lift over the WildGuard baseline of $0.075$,
while $d_l$-ablation provides only $+0.090$.

\subsection{Orthogonal decomposition: $d_l$'s effect is its $r$-component}
\label{sec:phase1-ortho}

We decompose $d_l = \alpha\,\hat{r}_l + d_l^\perp$ at $l^*_d = 26$,
then ablate each unit-normalised component and measure HarmBench
compliance under WildGuard. At $l^*_d = 26$: $\alpha = 32.80$,
$\|d_l\| = 82.80$, $\|d_l^\perp\| = 76.02$ ($91.8\%$ of $d_l$'s mass
is orthogonal to $r_l$), $\cos(d_l, r_l) = 0.396$.

\begin{table}[t]
\centering\small
\caption{Orthogonal decomposition at $l^*_d = 26$. By construction
$\hat{d}_\parallel = \mathrm{sign}(\alpha)\hat{r}$ when $\alpha > 0$,
so its row matches $\hat{r}$ exactly; $\hat{d}_\perp$ is the only row
with empirical content.}
\label{tab:phase1-ortho}
\begin{tabular}{l c c}
\toprule
Direction at $L26$  & val $S$  & HarmBench (WG) \\
\midrule
$\hat{r}$           & $+0.250$ & $0.635$        \\
$\hat{d}$           & $-0.031$ & $0.165$        \\
$\hat{d}_\parallel$ & $+0.250$ & $0.635$        \\
$\hat{d}_\perp$     & $\phantom{+}0.000$ & \textbf{0.005} \\
\bottomrule
\end{tabular}
\end{table}

The orthogonal complement $\hat{d}_\perp$ produces \emph{below-baseline}
compliance ($0.005 < 0.075$). The conclusion is mechanistic:
$d_l$ carries no refusal information independent of its $r$-aligned
component. Ablating $\hat{d}$ removes only $\cos(d, r) = 39.6\%$ of
the per-prompt $r$-projection that ablating $\hat{r}$ would remove,
which predicts the empirical compliance ($0.165$) sitting between
baseline and $\hat{r}$-ablation; the $91.8\%$ of $d$'s norm
orthogonal to $r$ does no causal work. The two extraction recipes
are not competing candidates for the refusal mediator: one identifies
it ($r_l$), and the other recovers a scaled version of the same
direction plus orthogonal capability-related noise ($d_l$).

% =====================================================================
\section{Phase 2: characterising EM geometry with the validated method}
\label{sec:phase2}

Having established that the within-class contrast captures the
refusal mediator, we apply the same method to FullFT-Medical to ask:
how does emergent-misalignment fine-tuning restructure the
residual-stream direction that mediates whatever refusal behaviour
remains?

\subsection{Cross-model setup and prompt re-filtering}
\label{sec:p2-setup}

Re-judging the same harmful candidate pool on EM and intersecting
with Phase 1's filter results yields $N_{\text{em}} = 242$ prompts
that both models refuse (out of $250$ candidates). Notable: $0$
prompts shifted from refused-by-instruct to compliant-on-EM, and EM
refused $7$ prompts instruct accepted. EM has not lost refusal
\emph{behaviour} on out-of-domain harmful prompts --- behavioural
preservation is essentially complete on AdvBench-style content. The
intersection size comfortably exceeds the $N{=}32$ floor demonstrated
by \citet{arditi2024refusal}, and we cap at the same $N{=}128$ used
in Phase 1 for direct comparability.

\subsection{Cross-model geometric agreement: a three-regime structure}
\label{sec:p2-geom}

For each layer $l$ we compute (a) the cross-model cosine
$\cos(r_{\text{em}}(l),\, \hat{r}_{l,\text{inst}}(l))$, (b) the
orthogonal decomposition
$r_{\text{em}}(l) = \beta(l)\,\hat{r}_{l,\text{inst}}(l)
                  + r_{\text{em}}^{\perp}(l)$
with $\beta(l) = \langle r_{\text{em}}(l),
\hat{r}_{l,\text{inst}}(l) \rangle$, and (c) the magnitude ratio
$\|r_{\text{em}}^\perp(l)\| / \|r_{\text{em}}(l)\|$.

\begin{figure}[t]
\centering
\fbox{\parbox{0.93\linewidth}{
  \centering\textcolor{red!85!black}{\textbf{[FIGURE 1 --- TODO]}}\\
  \small\textit{Cross-model geometric agreement across $48$ layers.
  Panel (a): $\cos(r_{\text{em}}, \hat{r}_{l,\text{inst}})$ per
  layer, with three regimes shaded. Panel (b):
  $\|r_{\text{em}}^\perp\| / \|r_{\text{em}}\|$ per layer.
  Vertical lines mark $l^*_r = 30$ and $l^*_{\text{em}} = 46$.
  Generated from \texttt{geometry\_em.json}.}
}}
\caption{Cross-model geometry across layers.}
\label{fig:p2-geom}
\end{figure}

The cosine curve breaks cleanly into three regimes:

\textbf{$L1$--$L18$ (shared)}. $\cos > 0.97$, perp fraction $< 0.22$.
Pre-refusal computation; both models look like the shared pretrained
base here. EM training barely modifies these layers.

\textbf{$L19$--$L28$ (divergence zone)}. Cosine drops $0.94 \to 0.43$
across ten layers, with the global minimum at $L26$
(cos $= 0.426$, perp\_frac $= 0.905$). Notably, this is exactly the
layer Phase 1 identified as $l^*_d$ --- the layer where
instruction-tuning concentrated its refusal-installing modification.
EM training has reorganised the representation most strongly at
exactly the layer instruct used to install it.

\textbf{$L29$--$L48$ (partial alignment)}. Cosine recovers to a
stable $0.66$--$0.74$ range; perp fraction stays around $0.68$--$0.78$.
At $l^*_{\text{em}} = 46$ specifically:
$\cos(r_{\text{em}}, \hat{r}_{l,\text{inst}}) = 0.692$ and
$\|r_{\text{em}}^\perp\|/\|r_{\text{em}}\| = 0.722$.

The mean cosine across all layers is $0.801$, the mean perp fraction
is $0.498$. Aggregate band IoU between the effective bands of $r_l$
(Phase 1) and $r_{\text{em}}$ (Phase 2) is $0.182$.

The within-EM geometry $\cos(r_{\text{em}}, d_{\text{em}})$ has mean
$-0.003$ across layers, ranging only $[-0.165, +0.140]$. Compared
with Phase 1's $\cos(r_l, d_l) = 0.589$ at $l^*_r$, the EM training
delta (instruct $\to$ EM on harmful prompts) is essentially
\emph{uncorrelated} with EM's within-class refusal contrast at every
layer.

\subsection{Sweep on EM: a 16-layer downstream shift and bimodal mediation}
\label{sec:p2-sweep}

Sweeping $r_{\text{em}}$ across all $48$ layers on the EM model
identifies $l^*_{\text{em}} = 46$ with peak $S = 0.594$ --- a
$16$-layer downstream shift from instruct's $l^*_r = 30$ and a
$0.406$-point reduction in single-direction ablation effect.
The effective band of $r_{\text{em}}$ contains only
$\{44, 45, 46, 47\}$ ($4$ layers, vs.\ $22$ for $r_l$); the IoU is
$0.182$.

The sweep curve is bimodal (Figure~\ref{fig:p2-sweep}): in addition
to the primary peak at $L44$--$L47$, a secondary band at $L29$--$L37$
shows non-trivial compliance ($S$ ranging $0.031$--$0.250$, peaking
at $L32$) --- exactly the layer range Phase 1 identified as $r_l$'s
broad effective band. We interpret this as a residual of the
original instruct refusal direction concentrated at $L30$--$L37$,
downgraded by EM training but not eliminated, plus a
newly-installed direction at the deeper $L44$--$L47$ region. The
random-direction control is $S = 0$ at every layer, confirming the
noise floor and that the bimodal structure is real causal signal.

\begin{figure}[t]
\centering
\fbox{\parbox{0.93\linewidth}{
  \centering\textcolor{red!85!black}{\textbf{[FIGURE 2 --- TODO]}}\\
  \small\textit{Per-layer $S_l$ for $r_{\text{em}}$ on EM (this
  work) and $r_l$ on instruct (Phase 1) on the same axes. Random
  control included for noise floor. Markers: $l^*_r = 30$, peak
  $S=1.000$; $l^*_{\text{em}} = 46$, peak $S=0.594$. Bimodal
  structure of $r_{\text{em}}$: secondary band $L29$--$L37$
  (peak $S=0.250$ at $L32$), primary band $L44$--$L47$.}
}}
\caption{Causal effect of single-layer ablation, EM (this work)
versus instruct (Phase 1).}
\label{fig:p2-sweep}
\end{figure}

\subsection{Three-way causal decomposition at $l^*_{\text{em}}$}
\label{sec:p2-3way}

To disentangle whether the residual mediation at $L46$ comes from
the instruct-aligned component of $r_{\text{em}}$ or from its
orthogonal addition, we ablate three unit-normalised directions on
the EM model at $L46$ and measure HarmBench compliance under both
the pattern judge and WildGuard. We include a random unit-vector
control, a baseline (no ablation), and a sanity row
$\hat{r}_{\text{em}}^{\parallel}
\equiv \mathrm{sign}(\beta)\,\hat{r}_{l,\text{inst}}$, whose ablation
must match $\hat{r}_{l,\text{inst}}$ exactly by construction
(Section~\ref{sec:phase1-ortho}).

The decomposition at $L46$: $\beta = 199.88$,
$\|r_{\text{em}}\| = 288.78$, $\|r_{\text{em}}^\parallel\| = 199.88$,
$\|r_{\text{em}}^\perp\| = 208.43$. The orthogonality check
$\langle r_{\text{em}}^\perp, \hat{r}_{l,\text{inst}}\rangle$ is
bounded by $1.07 \times 10^{-5}$, well within float32 noise.

\begin{table*}[t]
\centering\small
\caption{Three-way cross-model causal decomposition on EM at
$l^*_{\text{em}} = 46$. HarmBench compliance reported under both the
pattern judge and WildGuard; $\Delta$ is WildGuard minus pattern.
Capability is on $100$ Alpaca prompts (WildGuard).
$\hat{r}_{\text{em}}^{\parallel}$ matches $\hat{r}_{l,\text{inst}}$
exactly across all columns by construction.
$\hat{r}_{\text{em}}^{\perp}$ is statistically indistinguishable
from baseline and random under \emph{both} judges
($\Delta = 0$): the orthogonal complement carries no
refusal-mediating information.}
\label{tab:p2-3way}
\begin{tabular}{l c c c c c}
\toprule
Condition (at $L46$)
                                   & val $S$ (WG)
                                                & Pattern HB
                                                          & WildGuard HB
                                                                       & $\Delta$
                                                                                  & Capability \\
\midrule
baseline                           & $0.000$    & $0.025$ & $0.005$    & $-0.020$ & $0.970$    \\
ablate $\hat{r}_{\text{em}}$       & $\mathbf{0.594}$
                                                & $0.895$ & $\mathbf{0.495}$
                                                                       & $\mathbf{-0.400}$
                                                                                  & $0.990$    \\
ablate $\hat{r}_{l,\text{inst}}$   & $0.062$    & $0.560$ & $0.205$    & $-0.355$ & $0.990$    \\
ablate $\hat{r}_{\text{em}}^{\parallel}$ (sanity)
                                   & $0.062$    & $0.560$ & $0.205$    & $-0.355$ & $0.990$    \\
ablate $\hat{r}_{\text{em}}^{\perp}$
                                   & $0.000$    & $\mathbf{0.005}$
                                                          & $\mathbf{0.005}$
                                                                       & $\mathbf{\phantom{-}0.000}$
                                                                                  & $0.970$    \\
ablate random                      & $0.000$    & $0.035$ & $0.005$    & $-0.030$ & $0.960$    \\
\bottomrule
\end{tabular}
\end{table*}

\paragraph{Hypothesis interpretation.}
The three-way table distinguishes:
\begin{itemize}\itemsep=0pt
\item \textbf{(A) Dilution.} $S(\hat{r}_{\text{em}})$ substantial,
      $S(\hat{r}_{\text{em}}^\perp) \approx 0$. EM still mediates
      refusal through the instruct-aligned direction; the orthogonal
      addition is non-refusal residual.
\item \textbf{(B) Substitution.} $S(\hat{r}_{\text{em}})$
      substantial, $S(\hat{r}_{l,\text{inst}})$ small,
      $S(\hat{r}_{\text{em}}^\perp)$ comparable to
      $S(\hat{r}_{\text{em}})$. EM training has built a new
      direction orthogonal to instruct's; the original direction no
      longer mediates.
\item \textbf{(C) Hybrid.} Both old and new components carry
      partial mediation.
\end{itemize}

\paragraph{The data unambiguously support hypothesis A.}
Ablating $\hat{r}_{\text{em}}^\perp$ produces HarmBench compliance
$0.005$ under both judges --- identical to baseline ($0.005$
WildGuard, $0.025$ pattern), to random ($0.005$, $0.035$), and
$\Delta = 0$ between the two judges. The pattern judge, which
otherwise inflates compliance estimates by missing sophisticated
and multilingual refusals, finds no compliance to inflate. The
$72\%$ of $r_{\text{em}}$'s magnitude orthogonal to instruct's
direction does no causal work. The instruct-aligned component is
the operative one: $\hat{r}_{\text{em}}^\parallel$-ablation matches
$\hat{r}_{l,\text{inst}}$-ablation exactly across both judges
($0.560$ pattern, $0.205$ WildGuard). The full
$\hat{r}_{\text{em}}$-ablation produces a larger effect ($0.495$
WG) than $\hat{r}_{l,\text{inst}}$-ablation alone ($0.205$) because
$\hat{r}_{\text{em}}$ is the layer-46-specific optimal projection;
but the additional effect is not attributable to the orthogonal
component, which is independently inert. Capability is preserved
across all conditions ($0.96$--$0.99$).

\paragraph{The Phase 1--Phase 2 mirror.}
Phase 1 found $\hat{d}_l^\perp$-ablation produced HarmBench
compliance $0.005$ (Table~\ref{tab:phase1-ortho}). Phase 2 finds
$\hat{r}_{\text{em}}^\perp$-ablation produces $0.005$ under both
judges. These are the same number, obtained from structurally
different decompositions across different model paradigms. The
unified result across both phases is: \emph{whichever way a
candidate direction is produced --- post-training delta on instruct,
within-class contrast on EM --- decomposing it against instruct's
$r$ and ablating the orthogonal complement gives no causal effect.}
The within-class refusal contrast in instruct $r_l$ is the canonical
refusal mediator across both training paradigms tested here.

\paragraph{Multilingual fallback in Phase 2.}
The $\Delta = -0.400$ on $\hat{r}_{\text{em}}$-ablation is the
largest pattern-vs-WildGuard discrepancy in either phase. Of the
$200$ HarmBench prompts the pattern judge calls compliant under
$\hat{r}_{\text{em}}$-ablation, $80$ are sophisticated or
non-English refusals that the regex misses --- a $40\%$ false
compliance rate, larger than Phase 1's $\hat{d}$-ablation
discrepancy. The same artefact reappears in a different paradigm,
strengthening the methodological case for multilingual judging.

% =====================================================================
\section{Discussion}
\label{sec:discussion}

\paragraph{The within-class contrast as the right tool.}
Both Phase 1 and Phase 2 reach the same methodological conclusion:
the within-class refusal contrast captures the causal mediator. In
Phase 1 the alternative ($d_l$) recovers a scaled version of the
same direction plus orthogonal noise, with the orthogonal component
producing $0.005$ HarmBench compliance under WildGuard. In Phase 2
the candidate within-EM contrast ($r_{\text{em}}$) decomposes
analogously: its component along $\hat{r}_{l,\text{inst}}$
mediates refusal, while its orthogonal complement also produces
$0.005$ compliance under both judges. The exact numerical match of
the orthogonal-row HarmBench compliance across
Tables~\ref{tab:phase1-ortho} and~\ref{tab:p2-3way} is striking and
supports a unified picture: across both training paradigms tested
here, no information beyond the instruct-anchored within-class
refusal direction is causally implicated in refusal mediation.

\paragraph{The L26 connection, reconsidered.}
A coincidence worth flagging: Phase 1's $l^*_d = 26$ and the
cross-model geometric divergence minimum at $L26$ in Phase 2 are the
same layer. Both findings point to $L26$ as a structurally
privileged location in this model family --- the layer where
instruction-tuning concentrates its modification, and where any
further fine-tuning (EM in our case) re-modifies most strongly.
Whether this is a property of Qwen2.5-14B specifically or a general
feature of post-trained transformers is left to future work.

\paragraph{What EM training actually did.}
Combining the geometric, sweep, and ablation results: EM training
did not install a new orthogonal refusal direction (perp-row inert
under both judges). It also did not preserve the original mechanism
intact (peak $S$ weakened from $1.000$ to $0.594$, layer migrated
16 downstream). Instead, it produced a layer-shifted, weakened
version of instruct's mechanism: the operative computation now lives
at $L46$ rather than $L30$, with reduced rank-1 character, but the
same fundamental direction. The new optimum $r_{\text{em}}$ at $L46$
is a layer-specific reorientation of instruct's refusal direction
with non-refusal medical-domain content stacked on; the medical
content is geometrically substantial but causally invisible.

\paragraph{Implications.}
Because hypothesis A holds, the conference-paper extension naturally
points to multi-EM characterisation across the LoRA-Sport,
Steering-Medical, and FullFT-Medical span (which covers four orders
of magnitude in $\|d_{\text{em}}\|$) rather than to inverse-problem
characterisation of a hypothetical new direction. The structural
result here suggests that all three EM intensities should preserve
hypothesis A: the magnitude of EM-training perturbation should
affect how much the operative layer migrates and how much rank-1
strength is preserved, not which direction mediates.

% =====================================================================
\section{Limitations and conclusion}
\label{sec:lim}

The empirical comparison rests on a single seed for direction
extraction, two judges (pattern + WildGuard), one architecture
(Qwen2.5-14B), and one EM model (FullFT-Medical) for the Phase 2
applied study. Multi-seed extraction with bootstrap confidence
intervals on $r_l$ and $r_{\text{em}}$ would tighten the methodology
contribution; we leave this to the conference version. We defer
KL-divergence-based capability checks (cf.\
\citet{arditi2024refusal}) and report instead byte-identical Alpaca
completions under $\hat{r}$-ablation as a stronger but narrower
analogue. The methodology applies the chat template to the base
model when computing $d_l$, which is OOD for base; we make this
choice to preserve token-position alignment but acknowledge its
tradeoff.

We have systematically compared two refusal-direction extraction
recipes under a unified causal-ablation framework. Phase 1 shows the
within-class contrast captures the refusal mediator, while the
post-training delta is causally inert beyond its $r$-aligned
projection. Phase 2 applies the validated method to an
emergent-misalignment fine-tune and finds that EM training migrates
the operative refusal layer $16$ layers downstream
($L30 \to L46$), halves the rank-1 ablation effect ($S: 1.000 \to
0.594$), and produces a bimodal sweep curve indicating coexisting
mechanisms. The three-way decomposition at $L46$ shows that the
orthogonal-to-instruct component of $r_{\text{em}}$ is causally
inert under both judges (HarmBench compliance $0.005$,
$\Delta = 0$), confirming that EM training preserves rather than
substitutes the refusal-mediating direction --- it is a downstream
migration of the same mechanism, not a replacement.
Code, intermediate artefacts, and full completions are available at
\fillin{repo URL}.

% =====================================================================
\bibliographystyle{plainnat}
\bibliography{references}

\end{document}
